The over-identifying wedge is the part of the reform's price effect that silence escaped. Had silent listings received the same increase in attention as rated ones, their price relative to E-to-G listings would have moved by $\hat\beta_1\hat e^*$, as Section~\ref{sec:equilibrium} predicts; they received less, and the shortfall, priced at the cost of the silent pool, is $\hat\lambda$. A spillover of one means the reform made energy something buyers ask about whether or not a rating is shown; a spillover near zero means it made the rating something buyers read but not something they ask for. Part~(iii) states what the design cannot do. The two gaps pin down one number per market, the perceived cost of silence, and low attention with a modest default fits them as well as slightly higher attention with an even more optimistic one. Normalising the default to zero delivers the largest attention compatible with the data.\footnote{The same algebra in levels identifies relative attention at a point in time rather than the spillover. We do not estimate it, because across apartment markets the level gaps carry a large market-specific component that is unrelated to energy and that the differences-in-differences remove: the ratio of mean level gaps and the cross-market slope of one level gap on the other imply very different values, $\attPsiLevROMApt$ and $\attPsiLevSlopeApt$ post-reform. We use levels only for signs.}
